\begin{lemma}[The tested identity for the pseudoinverse of the smooth core]
Fix $n \geq 1$.  Let $f : \mathbb{R}^n \to \mathbb{R}$ be smooth with compact
support and with $f \in L^2(\gamma_n)$, and let $u \in L^2(\gamma_n)$ be the
class of $f$.  Let $\lambda$ denote Lebesgue measure on $\mathbb{R}^n$, write
$P_s$ for the Ornstein--Uhlenbeck semigroup of \noderef{OUSemigroup}, $\Pi_0$ for
the projection off the zeroth Hermite mode, and $w = L^{-1/2}u$ for the
pseudoinverse of \noderef{FracPow} at exponent $-1/2$.  Suppose given
\begin{enumerate}
  \item[(a)] a family $A : (0,\infty) \times \{1,\ldots,n\} \times \mathbb{R}^n
    \to \mathbb{R}$ with
    \[
      A_s(\partial_i f)(x)
        := A(s,i,x)
        = \int_{\mathbb{R}^n} \partial_i f\bigl(e^{-s}x
            + \sqrt{1-e^{-2s}}\,z\bigr)\,d\gamma_n(z) ,
    \]
    the Mehler average of $\partial_i f$ at time $s$;
  \item[(b)] a family $\Phi : \{1,\ldots,n\} \times \mathbb{R}^n \to \mathbb{R}$
    with
    \[
      \Phi_i(x) = \frac{1}{\sqrt\pi}\int_0^\infty s^{-1/2}
        \Bigl(e^{-s}\bigl(2^{-1/2}A_s(\partial_i f)(x)\bigr)\Bigr)\,ds ,
    \]
    its subordination integral;
  \item[(c)] a constant $N$ with $|A_s(\partial_i f)(x)| \leq N$ for all $s$, $i$
    and $x$;
  \item[(d)] joint continuity of $(s,x) \mapsto A_s(\partial_i f)(x)$ on
    $(0,\infty)\times\mathbb{R}^n$, for each $i$;
  \item[(e)] integrability of $s \mapsto s^{-1/2}e^{-s}$ on $(0,\infty)$; and
  \item[(f)] that every $\gamma_n$-almost-everywhere statement holds
    $\lambda$-almost everywhere.
\end{enumerate}
Then for every index $i$ and every smooth compactly supported
$\psi : \mathbb{R}^n \to \mathbb{R}$,
\[
  \int_{\mathbb{R}^n} w(x)\,\partial_i\psi(x)\,d\lambda(x)
    = -\int_{\mathbb{R}^n} \bigl(\sqrt2\,\Phi_i(x)\bigr)\,\psi(x)\,d\lambda(x) .
\]
\end{lemma}
\begin{proof}
Fix an index $i$ and a smooth compactly supported $\psi$, and set
\[
  \Lambda(v) = \int_{\mathbb{R}^n} v(x)\,\partial_i\psi(x)\,d\lambda(x),
  \qquad v \in L^2(\gamma_n) .
\]

\emph{Step 1: $\Lambda$ is a bounded linear functional on $L^2(\gamma_n)$.}
Let $K = \operatorname{supp}\psi$, a compact set.  The function
$\partial_i\psi$ is continuous, $\psi$ being smooth, and vanishes off $K$,
because a point outside $K$ has an open neighbourhood on which $\psi$ vanishes
identically and hence on which all its derivatives vanish.  So
$|\partial_i\psi|$ attains a maximum on $K$ by the extreme value theorem when
$K \neq \emptyset$, and $\partial_i\psi \equiv 0$ when $K = \emptyset$; in either
case $\|\partial_i\psi\|_\infty < \infty$.  By clause (3) of
\noderef{GaussL2LocallyIntegrable} applied to $K$ there is $C \geq 0$, depending
on $K$ but not on the function, such that every $g \in L^2(\gamma_n)$ is
Lebesgue integrable on $K$ with
$\int_K |g|\,d\lambda \leq C\,\|g\|_{L^2(\gamma_n)}$.  Hence
$v\,\partial_i\psi$ is $\lambda$-integrable and
\[
  |\Lambda(v)| \leq \|\partial_i\psi\|_\infty \int_K |v|\,d\lambda
    \leq \|\partial_i\psi\|_\infty\,C\,\|v\|_{L^2(\gamma_n)} .
\]
$\Lambda$ is linear by linearity of the integral, and well defined on
$L^2(\gamma_n)$ classes: two representatives agree $\gamma_n$-a.e., hence
$\lambda$-a.e.\ by (f), so the Lebesgue integral is unchanged.

\emph{Step 2: pushing $\Lambda$ inside the subordination integral.}
By clauses (1) and (2) of \noderef{SubordinationPseudoInverse} the
$L^2(\gamma_n)$-valued map $s \mapsto s^{-1/2}P_s\Pi_0u$ is Bochner integrable on
$(0,\infty)$ and $w = (\sqrt\pi)^{-1}\int_0^\infty s^{-1/2}P_s\Pi_0u\,ds$.
Apply mathlib's \texttt{ContinuousLinearMap.integral\_comp\_commSL}, which
commutes a continuous linear map with a Bochner integral, with
$X = (0,\infty)$ carrying Lebesgue
measure, with $E = L^2(\gamma_n)$ --- a real Banach space, being a Hilbert space
--- with $F(s) = s^{-1/2}P_s\Pi_0u$, Bochner integrable by the clause just
quoted, and with the continuous linear functional $\Lambda$ of Step 1.  It gives
$\int_0^\infty \Lambda(s^{-1/2}P_s\Pi_0u)\,ds
  = \Lambda\bigl(\int_0^\infty s^{-1/2}P_s\Pi_0u\,ds\bigr)$.  Multiplying by
$(\sqrt\pi)^{-1}$ and using linearity of $\Lambda$ both to pull that constant out
and to pull the scalar $s^{-1/2}$ out of $\Lambda(s^{-1/2}P_s\Pi_0u)$,
\begin{equation*}
  \Lambda(w) = \frac{1}{\sqrt\pi}\int_0^\infty s^{-1/2}\,\Lambda(P_s\Pi_0u)\,ds .
  \tag{1}
\end{equation*}

\emph{Step 3: the projection may be dropped inside $\Lambda$.}
Let $c = \langle u, h_0\rangle$ be the zeroth Hermite coefficient of $u$ with
respect to \noderef{HermiteBasis}.  By \noderef{FracPow} at exponent $0$, the
element $\Pi_0u$ has coefficients $0$ at $\alpha = 0$ and $u_\alpha$ at
$\alpha \neq 0$; the element $u - c\,h_0$ has, by linearity of the coefficient
map and orthonormality of the basis, coefficients $u_0 - c = 0$ at $\alpha = 0$
and $u_\alpha$ at $\alpha \neq 0$.  The coefficient map is injective, so
$\Pi_0 u = u - c\,h_0$.  Applying \noderef{OUSemigroupSpectral} to both sides,
the coefficients of $P_s\Pi_0u$ are $0$ at $\alpha = 0$ and
$e^{-s|\alpha|}u_\alpha$ at $\alpha \neq 0$, and those of $P_su - c\,h_0$ are
$u_0 - c = 0$ at $\alpha = 0$ and $e^{-s|\alpha|}u_\alpha$ at $\alpha \neq 0$; by
injectivity again $P_s\Pi_0 u = P_s u - c\,h_0$.

By \noderef{HermiteL2ZeroAEOne} the class $h_0$ agrees $\gamma_n$-a.e.\ with the
constant function $1$, hence $\lambda$-a.e.\ by (f), so
$\Lambda(c\,h_0) = c\int_{\mathbb{R}^n}\partial_i\psi\,d\lambda$.  That integral
vanishes, by \noderef{CompIntegrationByParts} at the index $j := i$ with the
following data.  Take $q = \mathrm{id}_{\mathbb{R}}$ and $q' \equiv 1$: then $q$
is differentiable at every $t$ with derivative $1$, and $q'$, being constant, is
continuous.  Take the inner function to be $\equiv 1$, the constant function on
$\mathbb{R}^n$: it is differentiable everywhere with vanishing derivative, so its
$i$th partial is the zero function, which is continuous.  Take the test function
to be $\psi$.  The conclusion of that node reads
$\int q(1)\,\partial_i\psi\,d\lambda
  = -\int \bigl(q'(1)\cdot 0\bigr)\psi\,d\lambda$, i.e.
$\int\partial_i\psi\,d\lambda = 0$.  Hence $\Lambda(c\,h_0) = 0$ and
\begin{equation*}
  \Lambda(P_s\Pi_0u) = \Lambda(P_su) \qquad\text{for every } s>0. \tag{2}
\end{equation*}

\emph{Step 4: the identity at a fixed time.}
Fix $s > 0$.  Since $u$ is the class of $f$ and, by \noderef{MehlerFixedTimeTesting},
the coercion of $P_su$ agrees $\gamma_n$-a.e.\ with $A_sf$ --- hence
$\lambda$-a.e.\ by (f) --- we get
$\Lambda(P_su) = \int (A_sf)\,\partial_i\psi\,d\lambda$.  Apply
\noderef{SemigroupGradientCommute} at this $s$, to this $f$, and at the index
$i$: its clause (4) says $A_sf$ is differentiable at every $x$ with
$\partial_i(A_sf)(x) = e^{-s}A_s(\partial_i f)(x)$, and its clause (3) says
$A_s(\partial_i f)$ is continuous.  Hence $x \mapsto \partial_i(A_sf)(x)$, the
constant $e^{-s}$ times a continuous function, is continuous.

These are exactly the two hypotheses \noderef{CompIntegrationByParts} places on
its inner function at the index $j := i$: differentiability everywhere, and
continuity of the single partial $\partial_i$.  Apply it with inner function
$A_sf$, with $q = \mathrm{id}_{\mathbb{R}}$, $q' \equiv 1$, and test function
$\psi$.  Its conclusion reads
\[
  \int (A_sf)\,\partial_i\psi\,d\lambda
    = -\int \partial_i(A_sf)\,\psi\,d\lambda
    = -e^{-s}\int A_s(\partial_i f)\,\psi\,d\lambda ,
\]
the last equality by clause (4).  Convergence of both integrals is part of that
node's conclusion, and there is no boundary term because $\psi$ is compactly
supported.  Combining with $(2)$,
$\Lambda(P_s\Pi_0u) = -e^{-s}\int A_s(\partial_i f)\,\psi\,d\lambda$.

\emph{Step 5: Fubini in $(s,x)$.}
Substituting the last display into $(1)$,
\[
  \Lambda(w) = -\frac{1}{\sqrt\pi}\int_0^\infty\!\!\int_{\mathbb{R}^n}
      s^{-1/2}e^{-s}\,A_s(\partial_i f)(x)\,\psi(x)\,d\lambda(x)\,ds .
\]
The integrand is jointly measurable: $(s,x) \mapsto A_s(\partial_i f)(x)$ is
jointly continuous by (d), and $(s,x) \mapsto s^{-1/2}e^{-s}\psi(x)$ is
continuous.  It is absolutely integrable, since by (c), the vanishing of $\psi$
off the compact set $K$, the finiteness of $\lambda(K)$ and (e),
\[
  \int_0^\infty\!\!\int
      \bigl|s^{-1/2}e^{-s}A_s(\partial_i f)(x)\,\psi(x)\bigr|\,d\lambda(x)\,ds
    \leq N\,\|\psi\|_\infty\,\lambda(K)\int_0^\infty s^{-1/2}e^{-s}\,ds
    < \infty .
\]
Both factor measures are $\sigma$-finite, so Fubini permits exchanging the order:
\[
  \Lambda(w) = -\int_{\mathbb{R}^n}\psi(x)
      \Bigl[\frac{1}{\sqrt\pi}\int_0^\infty s^{-1/2}e^{-s}
        A_s(\partial_i f)(x)\,ds\Bigr]d\lambda(x)
    = -\int_{\mathbb{R}^n} \sqrt2\,\Phi_i(x)\,\psi(x)\,d\lambda(x) ,
\]
the last equality because, by (b), $\Phi_i(x)$ is
$(\sqrt\pi)^{-1}\int_0^\infty s^{-1/2}e^{-s}2^{-1/2}A_s(\partial_i f)(x)\,ds$, so
multiplying by $\sqrt2$ removes the factor $2^{-1/2}$.  Since
$\Lambda(w) = \int w\,\partial_i\psi\,d\lambda$, this is the stated identity.
\end{proof}
